\section{Limitations}
The evidence concerns two closed-set tasks, one checkpoint, three seeds, and one implemented fixed-head control. Each task has a fixed template and candidate set; no unseen-task, paraphrase, or new-label evaluation was performed. Candidate scoring's measured advantage cannot be assigned solely to textual label semantics because the control changes pooling and input structure. Specialized biological encoders and generative sequence--text systems were not evaluated under a matched budget.

Data independence remains limited. Exact and DNA reverse-complement overlap checks do not establish homology-separated protein evaluation or species-separated DNA evaluation. Historical cluster summaries lack a verified membership mapping to these local records. The biological BPE was fitted on external corpora whose overlap with the benchmark is not fully established, and the parent checkpoint's historical exposure is not excluded. The inherited eligibility filter also restricts coverage. These limitations rule out an unqualified leakage-free generalization claim.

This was an iterative development study, not a publicly preregistered experiment. Earlier compact-vocabulary runs informed development; the final twelve-model protocol was fixed before its test predictions. Seed SDs do not quantify biological sampling uncertainty, and no significance claim is made. The observed order sensitivity and incomplete probability calibration limit direct deployment claims. The study is a computational benchmark, with no experimental or clinical validation.
